\documentclass[a4paper]{article}

\usepackage[fontset=fandol]{ctex}
\usepackage{amsmath, amssymb}
\usepackage{graphicx}
\usepackage[margin=2.5cm]{geometry}
\usepackage[most]{tcolorbox}
\usepackage{etoolbox}
\usepackage{listings}
\usepackage{hyperref}
\usepackage{booktabs}
\usepackage{subcaption}
\usepackage{float}
\usepackage{tikz}
\IfFileExists{pgfplots.sty}{
    \usepackage{pgfplots}
    \pgfplotsset{compat=1.18}
}{}

\newtcolorbox{knowledgebox}[1]{
    enhanced, colback=blue!5!white, colframe=blue!75!black, colbacktitle=blue!75!black,
    coltitle=white, fonttitle=\bfseries, title=#1, attach boxed title to top left={yshift=-2mm, xshift=2mm},
    boxrule=1pt, sharp corners
}

\newtcolorbox{importantbox}[1]{
    enhanced, colback=yellow!10!white, colframe=yellow!80!black, colbacktitle=yellow!80!black,
    coltitle=black, fonttitle=\bfseries, title=#1, sharp corners
}

\newtcolorbox{warningbox}[1]{
    enhanced, colback=red!5!white, colframe=red!75!black, colbacktitle=red!75!black,
    coltitle=white, fonttitle=\bfseries, title=#1, sharp corners
}

\lstset{
    language=Python,
    basicstyle=\ttfamily\small,
    keywordstyle=\color{blue},
    stringstyle=\color{red!60!black},
    commentstyle=\color{green!60!black},
    breaklines=true,
    frame=single,
    numbers=left,
    numberstyle=\tiny\color{gray},
    captionpos=b,
    extendedchars=false
}

\newcommand{\notetitle}{CS224R Lecture 11: Model-Based RL}
\newcommand{\noteauthors}{Codex}
\newcommand{\notedate}{2026-04-13}
\newcommand{\videochannel}{Stanford CS224R}
\newcommand{\videopublishdate}{2025-05-07}
\newcommand{\videoduration}{约 01:13:20}
\newcommand{\videourl}{https://www.youtube.com/watch?v=PvqyGnOirgA}
\newcommand{\videocoverpath}{cover.jpg}

\begin{document}

\begin{titlepage}
\centering
{\Large 课程笔记\par}
\vspace{1.2cm}
{\huge\bfseries \notetitle\par}
\vspace{0.8cm}
{\large \noteauthors\par}
\vspace{0.3cm}
{\large \notedate\par}
\vspace{1.2cm}

\includegraphics[width=0.82\textwidth,height=0.45\textheight,keepaspectratio]{\videocoverpath}\par

\vfill
\begin{tcolorbox}[width=0.9\textwidth, colback=black!2!white, colframe=black!60, sharp corners]
\textbf{视频作者/频道}：\videochannel\par
\textbf{发布日期}：\videopublishdate\par
\textbf{视频时长}：\videoduration\par
\textbf{视频链接}：\href{\videourl}{\nolinkurl{\videourl}}
\end{tcolorbox}
\end{titlepage}

\tableofcontents
\newpage

\section{Model-Based RL 的核心问题：能不能先学一个世界模型}

这节课把强化学习分成一个非常清楚的谱系：offline / online、imitation / RL、model-free / model-based。Model-based RL 的核心出发点并不复杂，就是先学一个能预测未来的模型，再利用这个模型帮助决策。讲者把它称为 ``learned simulator''，这比把它仅仅理解成下一状态预测器更准确，因为它承担的是环境近似器的角色。

\begin{figure}[H]
\centering
\includegraphics[width=0.9\textwidth]{pdf_pages/page-06.png}
\caption{learned simulator 的适用场景：物理系统、视频预测、金融市场、游戏规则等，都可以视为状态转移规律的建模问题。}
\end{figure}

讲者特别强调，``模型'' 并不一定意味着高精度物理方程。对于机器人和驾驶，模型可能近似的是显式物理世界；对于难以观测底层变量的场景，模型更像 action-conditioned video prediction；在游戏中，规则本身就是环境模型；而在多智能体设置里，其他玩家有时也被当作动力学的一部分。

因此，model-based RL 的第一层判断不是 ``模型有没有用''，而是 ``这个领域有没有希望学出足够有用的模型''。如果答案是否定的，那么后面所有 planning 与 imagined rollouts 都会被系统性误导。

\subsection{一个最基础的 learned simulator pipeline}

课程给出的第一版 pipeline 非常朴素：

\begin{enumerate}
\item 用某个基础策略采集数据集 $\mathcal{D}=\{(s,a,s')_i\}$；
\item 训练动力学模型 $f_\phi(s,a)$ 去预测 $s'$；
\item 在学到的模型内部做动作优化或规划；
\item 将规划后的动作拿回真实环境执行。
\end{enumerate}

最直接的训练目标就是最小化一步预测误差：
\[
\min_\phi \sum_i \|f_\phi(s_i,a_i)-s_i'\|_2^2.
\]
这一定义本身已经暴露了 model-based RL 的两个关键前提：一是你必须有覆盖足够广的数据，二是你希望一步预测误差能够外推到多步决策。

\begin{knowledgebox}{本讲主线}
Lecture 11 的核心不是背算法名，而是建立一个判断框架：学到的模型如何被使用？为什么 planning 会失败？什么时候该 replan？什么时候应该直接用模型生成额外数据，再交给 model-free RL 来更新策略？
\end{knowledgebox}

\subsection{本章小结}

开场部分建立了 model-based RL 的最基本视角：先学模型，再用模型。真正困难的地方不在于写下 $f_\phi(s,a)$，而在于这个模型到底要被怎样调用，以及它的误差会在决策中如何被放大。

\section{第一条路：在模型里直接规划}

\subsection{把动作序列当作优化变量}

如果模型 $f_\phi$ 已经能近似预测未来，那么一个自然做法是直接对未来动作序列做优化。对于规划视角，目标可以写成
\[
\max_{a_{t:t+H-1}} \sum_{h=0}^{H-1} r(s_{t+h}, a_{t+h})
\quad
\text{s.t.}\quad
s_{t+h+1}=f_\phi(s_{t+h},a_{t+h}).
\]

这里真正被优化的不是网络参数，而是未来的动作序列 $a_t,\dots,a_{t+H-1}$。这和常规神经网络训练很不一样：我们把模型固定下来，然后在模型内部搜索一串最优控制。

\begin{figure}[H]
\centering
\includegraphics[width=0.92\textwidth]{pdf_pages/page-12.png}
\caption{通过 learned model 直接对动作序列做规划：先学 $f_\phi(s,a)$，再在模型里搜索高回报动作。}
\end{figure}

讲者把这类做法称为 planning。它的优势是：不需要先学出一个显式 policy 网络，只要模型和奖励足够可信，就能在当前状态上 ``向前想''。但代价也很直接，推理时会很重，因为每到一个状态都得重新做优化。

\subsection{梯度法与采样法的差异}

课程把规划里的动作优化分成两大类。第一类是\textbf{梯度法}，即 reward 对动作可微时，通过反向传播直接更新动作序列。第二类是\textbf{采样法}，即不要求可微，只需要能评估某串动作的累计回报。

\begin{figure}[H]
\centering
\includegraphics[width=0.82\textwidth]{pdf_pages/page-13.png}
\caption{梯度优化与采样优化的对比：前者依赖可微结构，后者更黑盒，但通常更稳健。}
\end{figure}

讲者对采样法的说明尤其重要。最简单的是 random shooting：随机采样很多候选动作序列，逐个 rollout，选回报最高的一组。更高级的是 cross-entropy method, CEM：保留 elite sequences，再拟合一个新的采样分布迭代更新。CEM 的优点是简单、易并行、对不可微模型友好；缺点是动作维度和规划 horizon 一大，就会迅速变贵。

\begin{figure}[H]
\centering
\includegraphics[width=0.9\textwidth]{pdf_pages/page-16.png}
\caption{CEM 式采样规划：采样候选动作序列、评估回报、保留 elite、再拟合采样分布。}
\end{figure}

\subsection{本章小结}

如果模型可用，那么 planning 是最直接的 model-based RL 方案。它的思想简单、目标清晰，但它把困难从 ``学策略'' 转移到了 ``每一步都要在线搜索''。因此，planning 的实际可行性高度依赖模型是否可靠、动作空间是否可控，以及测试时预算是否允许。

\section{为什么直接规划会失败，以及 MPC 如何补救}

\subsection{数据分布失配会让模型在关键区域失真}

讲者用一个 ``向右走会上高地，但继续走会掉下悬崖'' 的例子解释 model-based planning 的经典失败模式。训练数据来自某个基础策略 $\pi_0$，而最终规划出的策略 $\pi_f$ 会访问完全不同的状态区域，于是出现
\[
p_{\pi_0}(s)\neq p_{\pi_f}(s).
\]
模型在训练数据覆盖的区域也许很准，但在新访问到的状态上会产生系统性误差。规划器又会主动寻找高回报区域，于是极容易被这些未见区域里的虚假高值吸引。

\begin{figure}[H]
\centering
\includegraphics[width=0.84\textwidth]{pdf_pages/page-18.png}
\caption{direct planning 的典型失败：模型只在旧策略数据分布上准确，到了新状态区域就会犯系统性错误。}
\end{figure}

这也是为什么 model-based RL 不能只离线学一次模型就完事。讲者在口头解释里明确说到：当你执行规划动作后，应把新数据加入数据集、重新拟合模型，再继续迭代。只有这样，模型才会逐渐覆盖最终策略真正访问的区域。

\subsection{open-loop 为什么脆弱}

另一类失败并不来自训练分布，而来自执行方式。若你一次规划出整串动作 $a_{t:t+H-1}$，然后不再看环境反馈、直接全都执行完，这就是 open-loop control。它的问题是：任何一个中途偏差，无论来自环境扰动还是模型误差，都会让后续动作建立在错误假设上。

\begin{figure}[H]
\centering
\includegraphics[width=0.9\textwidth]{pdf_pages/page-22.png}
\caption{MPC 的核心思想是 ``plan \& replan''：每次只执行第一个动作，再根据真实新状态重新规划。}
\end{figure}

讲者在课堂讨论里给了一个非常工程化的解释：如果世界变化很快，或者模型没法精确预测未来，那么 open-loop 的长动作序列其实相当危险，因为你根本没有机会根据真实反馈修正自己。

\subsection{MPC：planning 也要闭环}

为了解决 open-loop 的脆弱性，课程引入 model-predictive control, MPC。它仍然规划一个有限 horizon 的动作序列，但只执行第一个动作；得到真实下一状态后，再重新规划剩余序列。这相当于把 planning 放进一个闭环反馈系统里。

MPC 的优点很明确：

\begin{itemize}
\item 可以根据真实新状态修正偏差；
\item 对局部模型误差更鲁棒；
\item 能把长程任务分解成一系列短程滚动规划问题。
\end{itemize}

它的代价同样明显：每一步都要重新规划，所以测试时计算开销很大。讲者还特别指出，planning 通常只适合短 horizon 问题，因为一方面长 horizon 的搜索太贵，另一方面模型反复 rollout 的误差会累积。

\subsection{本章小结}

这一部分给出了 planning 能不能落地的判断标准。只要出现数据分布失配、环境变化、模型误差累积，单次 open-loop 规划都会变得不可靠。MPC 的意义不在于把 planning 变简单，而在于把它从脆弱的离线计划，变成一个能利用环境反馈持续纠错的闭环过程。

\section{第二条路：不用模型直接决策，而是用模型扩充数据}

\subsection{为什么 imagined rollouts 更实用}

讲者随后转向另一条思路：既然 planning 测试时代价大、长 horizon 下也不稳，那不如用 learned model 生成额外训练数据，再交给熟悉的 model-free RL 更新 policy。这就是 model-based policy optimization 的核心直觉。

关键不是从初始状态生成很长的虚拟轨迹，而是\textbf{从真实数据中的很多状态出发，生成短的 synthetic rollouts}。这样做能在真实数据附近增加局部覆盖，同时避免模型在长 horizon 下越滚越偏。

\begin{figure}[H]
\centering
\includegraphics[width=0.9\textwidth]{pdf_pages/page-28.png}
\caption{model-based policy optimization 的全算法：真实环境数据训练模型，模型再从真实状态附近生成短 imagined rollouts。}
\end{figure}

课程给出的完整 pipeline 是：

\begin{enumerate}
\item 用当前策略收集真实数据，加入 $D_{\text{env}}$；
\item 用 $D_{\text{env}}$ 更新动力学模型 $p_\theta(s'|s,a)$；
\item 从 $D_{\text{env}}$ 的状态出发，在模型里用当前策略采样短 synthetic rollouts，加入 $D_{\text{model}}$；
\item 用 $D_{\text{model}}$ 更新策略和 critic，必要时也可混入 $D_{\text{env}}$。
\end{enumerate}

这里最重要的工程经验是 ``short rollouts from many real states''。讲者在口头解释里强调过，从真实轨迹末尾往后接长 rollout 并不稳妥；更好的做法是从 replay buffer 里很多中间状态出发，生成少量、短程、局部可信的 imagined transitions。

\subsection{什么时候 model-based 值得做}

课程没有把 model-based RL 神化，而是很克制地列出了利弊。优点包括：

\begin{itemize}
\item 如果模型好学，样本效率会显著提高；
\item 模型可以在没有 reward 标注的情形下，用自监督方式训练；
\item 模型有时是 task-agnostic 的，可以跨任务复用。
\end{itemize}

缺点同样扎实：

\begin{itemize}
\item 模型并不直接为任务回报优化；
\item 有些领域里学模型比学策略还难；
\item 额外引入模型，意味着额外超参数、额外算力、额外调试面。
\end{itemize}

\begin{warningbox}{课程给出的判断标准}
是否使用 model-based RL，不该由 ``它看起来更聪明'' 决定，而应由 ``这个领域中的动力学是否足够可学'' 决定。模型越难学，planning 与 imagined rollouts 的收益就越不稳定。
\end{warningbox}

\subsection{本章小结}

model-based policy optimization 可以看成 planning 与 model-free RL 的折中：不在测试时频繁搜索，而是在训练时用模型补充局部覆盖。它特别适合真实交互成本高、但局部动力学又相对可学的场景。

\section{案例：灵巧手操控为什么适合做 model-based RL}

\subsection{任务定义与为什么它是好案例}

讲者最后用 dexterous manipulation 作为案例，原因有两个。第一，五指灵巧手是高维、易碎、真实交互代价高的系统，样本效率特别关键；第二，这个任务能清楚展示 model-based 方法的关键部件，比如 ensemble、CEM planner、reward design 与真实系统的闭环执行。

\begin{figure}[H]
\centering
\includegraphics[width=0.86\textwidth]{pdf_pages/page-32.png}
\caption{dexterous manipulation 被当作 case study 的原因：高维、真实交互昂贵，但又迫切需要高样本效率。}
\end{figure}

在这个案例里，状态包括手和物体的位置；动作空间是控制 24 自由度左右的五指手；奖励目标是跟踪指定物体轨迹，并惩罚掉落。这种任务如果完全依赖 model-free RL，真实世界数据成本会很高。

\subsection{为什么 ensemble 和改造版 CEM 有帮助}

课程提到，案例中使用的模型并不是单个网络，而是一个 ensemble。原因在于，不同模型会在不同区域犯不同错误，集成后能减轻单一模型过度乐观的问题。规划器方面，则采用 CEM 的变体，并加上 smoother action sequence 等技巧，以便搜索更平滑、可执行的动作序列。

\begin{figure}[H]
\centering
\includegraphics[width=0.88\textwidth]{pdf_pages/page-34.png}
\caption{仿真实验结果表明：在该任务上，合适的 model-based 方法既比 model-free 更省样本，也比较弱的 model-based baseline 更强。}
\end{figure}

这恰好对应本讲前半段的理论：如果模型误差会误导规划，那 ensemble 就是一种减轻误差 exploitation 的办法；如果规划动作过于抖动，那就要在 planner 内加入更贴近物理执行的约束。

\subsection{真实世界部署说明了什么}

课程最后展示了 ShadowHand 上的真实系统结果：Baoding balls 旋转等任务可以在有限小时级别的数据内学成。这说明 model-based RL 在硬件脆弱、交互昂贵的场景里尤其有吸引力，因为它能把有限真实经验 ``放大'' 成更多有价值的决策信号。

\begin{figure}[H]
\centering
\includegraphics[width=0.86\textwidth]{pdf_pages/page-36.png}
\caption{真实世界 ShadowHand 实验凸显了 model-based RL 的核心价值：在昂贵硬件上用更少真实交互获得可用策略。}
\end{figure}

\subsection{本章小结}

案例部分把整节课的抽象讨论落到了具体系统上。高维硬件控制之所以适合 model-based RL，是因为真实数据太贵，而局部动力学又并非完全不可学。此时，ensemble model、CEM planner、MPC 式闭环执行就会组合成一套可落地的工程方案。

\section{总结与延伸}

Lecture 11 的关键贡献，是把 ``学模型有什么用'' 这个问题拆成了三层。第一层是世界模型本身：你能否从数据里学到足够可信的动力学近似。第二层是决策调用：你是直接在模型里做 planning，还是把模型当成 imagined data generator。第三层是误差管理：你怎样处理数据分布失配、open-loop 脆弱性、长 horizon 误差累积，以及规划器对模型漏洞的利用。

从课程视角看，planning 与 model-free RL 并不是对立面。planning 适合短 horizon、需要强在线反应的场景；model-based policy optimization 适合用 imagined rollouts 提高样本效率；MPC 则是让 planning 进入闭环控制的关键桥梁。所有这些方法的共同前提都是：模型至少在某个局部区域里足够有用。

如果把这节课和现代 RL、机器人系统联系起来，可以得到一个很务实的结论：model-based RL 不是默认更高级，而是默认更脆弱。只有在动力学确实可学、硬件交互成本确实高、而你又有足够工程能力管理模型误差时，它才会成为真正有优势的路线。这也是为什么课程最后会落在 dexterous manipulation 这样的案例上，因为它恰好满足了这些条件。

\end{document}
